% Resumo em portugues.
\setlength{\absparsep}{18pt}

\begin{resumo}
\textbf{Contexto:} Sistemas intensivos em dados são fundamentais para o processamento, armazenamento, movimentação e análise de grandes volumes de dados em infraestruturas modernas de software. Esses sistemas envolvem decisões arquiteturais relacionadas a paradigmas de processamento, ingestão de dados, escalabilidade, uso de recursos e tolerância a falhas, as quais impactam diretamente atributos como \textit{throughput}, latência, estabilidade e custo computacional. \textbf{Problema:} Apesar da existência de diferentes técnicas, \textit{frameworks} e estratégias de otimização para sistemas intensivos em dados, as evidências empíricas disponíveis na literatura permanecem fragmentadas, heterogêneas e, muitas vezes, difíceis de comparar. Essa limitação dificulta a tomada de decisão arquitetural baseada em evidências, especialmente em cenários que exigem a escolha entre processamento em lote, processamento em fluxo e diferentes configurações de execução. \textbf{Solução:} Esta pesquisa propõe investigar uma abordagem quantitativa de apoio à decisão arquitetural em sistemas intensivos em dados, integrando evidências da literatura, resultados experimentais e características observáveis dos dados, como volume, entropia, redundância, cardinalidade e variabilidade. \textbf{Método:} A pesquisa foi organizada em três etapas complementares: uma Revisão Sistemática da Literatura, uma campanha experimental controlada e a proposta de um modelo quantitativo de apoio à decisão arquitetural. A revisão sistemática identificou técnicas, métricas, lacunas e evidências relacionadas ao design de sistemas intensivos em dados. A campanha experimental comparou processamento \textit{batch} e \textit{streaming} em um ambiente baseado em Apache Spark, Spark Structured Streaming, Apache Kafka e Docker, utilizando o dataset NYC Taxi. Foram avaliados diferentes volumes de dados, taxas de eventos, números de \textit{workers} e intervalos de \textit{trigger}, considerando métricas como \textit{throughput}, latência, uso de CPU, consumo de memória, \textit{backpressure} e escalabilidade. \textbf{Sumarização dos resultados:} A revisão sistemática identificou 979 registros, dos quais dez estudos primários atenderam aos critérios de inclusão e qualidade. A síntese quantitativa exploratória indicou efeitos positivos de otimizações arquiteturais, com razão média de resposta de aproximadamente $2{,}7\times$ em relação às configurações de referência. Os resultados experimentais parciais indicaram que o processamento \textit{batch} apresentou maior \textit{throughput} nos cenários avaliados, enquanto o processamento \textit{streaming} mostrou maior sensibilidade à taxa de eventos, ao intervalo de \textit{trigger} e aos custos de coordenação dos micro-lotes. \textbf{Contribuição:} Espera-se que esta pesquisa contribua para a tomada de decisão arquitetural em sistemas intensivos em dados por meio de uma abordagem baseada em evidências empíricas, resultados experimentais e características informacionais dos dados, auxiliando arquitetos de software e engenheiros de dados na comparação preliminar entre alternativas arquiteturais antes da implementação completa de uma solução.

\textbf{Palavras-chave:} Sistemas Intensivos em Dados. Arquitetura de Software. Tomada de Decisão Arquitetural. Processamento em Fluxo. Teoria da Informação.


\end{resumo}
